% GENERATED FILE. Edit canonical lessons, then run npm run generate:books.
% canonical-source-hash: fnv1a64:660112849e226152
% canonical-lessons: JA-C154-fudan, JA-C154-kuse, JA-C154-nichijou, JA-C154-kurasu, JA-C154-nareru

\chapter{Habits and Everyday Life}
\label{ch:ja-habits-and-everyday-life}
\hlchaptermodality{154}

\begin{chapteropening}
\textbf{By the end of this chapter:} \emph{I can describe my usual days and habits.}

Complete the last lesson of chapter 154: I can describe my usual days and habits.
\end{chapteropening}

\section[\texorpdfstring{\ja{ふだん} (fudan)}{ふだん (fudan)}]{\ja{ふだん} (fudan) — usually; everyday}

\label{lesson:JA-C154-fudan}



\begin{quote}
Before the new one: say the Japanese for to lose, then the Japanese for to be in time.
\end{quote}

\subsection*{You'll want to know: \ja{ふだん}}
\textbf{\ja{ふだん}} — \emph{fudan} — \textquotedblleft{}usually; everyday\textquotedblright{}.

\begin{quote}
\textbf{\ja{ふだんは} \ja{うちで} \ja{ごはんを} \ja{たべます。}} — \emph{fudan wa uchi de gohan o tabemasu} — \textquotedblleft{}I usually eat at home.\textquotedblright{}
\end{quote}

\begin{cousinweb}
\textbf{\ja{ふだん}} (\emph{fudan}) is Sino-Japanese (kango): two Chinese characters read with their on'yomi, readings borrowed from Chinese. They are \emph{fu}, \textquotedblleft{}general, widely\textquotedblright{} (Middle Chinese \emph{phuX}) and \emph{dan}, \textquotedblleft{}step, grade\textquotedblright{} (\emph{dwanH}) (KANJIDIC2; Middle Chinese in Baxter's transcription).
\end{cousinweb}

\begin{grammarlens}[title={What you've built}]
One more word for everyday life at A2.
\end{grammarlens}

\subsection*{Your turn}
\begin{itemize}
\raggedright
  \item \emph{Say it:} \emph{fudan}
  \item \emph{Say it:} \emph{fudan}, once more
  \item \emph{Say it:} \emph{fudan wa uchi de gohan o tabemasu} — I usually eat at home.
  \item \emph{Recall:} say the Japanese for to lose, then the Japanese for to be in time, then say \emph{fudan} again
\end{itemize}

\begin{tcolorbox}[breakable,colback=teal!4,colframe=teal!35!black,title={Before you move on}]
What does \ja{ふだん} mean? (\textquotedblleft{}Usually; everyday\textquotedblright{}.) Say it once more.
\end{tcolorbox}
\section[\texorpdfstring{\ja{くせ} (kuse)}{くせ (kuse)}]{\ja{くせ} (kuse) — a habit; a quirk}

\label{lesson:JA-C154-kuse}



\begin{quote}
Before the new one: say the Japanese for to be in time, then the Japanese for usually.
\end{quote}

\subsection*{You'll want to know: \ja{くせ}}
\textbf{\ja{くせ}} — \emph{kuse} — \textquotedblleft{}a habit; a quirk\textquotedblright{}.

\begin{quote}
\textbf{\ja{それは} \ja{ちちの} \ja{くせです。}} — \emph{sore wa chichi no kuse desu} — \textquotedblleft{}That is my father's habit.\textquotedblright{}
\end{quote}

\begin{cousinweb}
\textbf{\ja{くせ}} (\emph{kuse}) is a native Japanese word (wago), inherited rather than borrowed from Chinese; its kanji was attached later as a kun'yomi, and its older origin is not settled (JMdict, KANJIDIC2).
\end{cousinweb}

\begin{grammarlens}[title={What you've built}]
One more word for everyday life at A2.
\end{grammarlens}

\subsection*{Your turn}
\begin{itemize}
\raggedright
  \item \emph{Say it:} \emph{kuse}
  \item \emph{Say it:} \emph{kuse}, once more
  \item \emph{Say it:} \emph{sore wa chichi no kuse desu} — That is my father's habit.
  \item \emph{Recall:} say the Japanese for to be in time, then the Japanese for usually, then say \emph{kuse} again
\end{itemize}

\begin{tcolorbox}[breakable,colback=teal!4,colframe=teal!35!black,title={Before you move on}]
What does \ja{くせ} mean? (\textquotedblleft{}A habit; a quirk\textquotedblright{}.) Say it once more.
\end{tcolorbox}
\section[\texorpdfstring{\ja{にちじょう} (nichijō)}{にちじょう (nichijō)}]{\ja{にちじょう} (nichijō) — everyday life; daily}

\label{lesson:JA-C154-nichijou}



\begin{quote}
Before the new one: say the Japanese for usually, then the Japanese for a habit.
\end{quote}

\subsection*{You'll want to know: \ja{にちじょう}}
\textbf{\ja{にちじょう}} — \emph{nichijō} — \textquotedblleft{}everyday life; daily\textquotedblright{}.

\begin{quote}
\textbf{\ja{にちじょうは} \ja{たのしいです。}} — \emph{nichijō wa tanoshii desu} — \textquotedblleft{}Everyday life is enjoyable.\textquotedblright{}
\end{quote}

\begin{cousinweb}
\textbf{\ja{にちじょう}} (\emph{nichijō}) is Sino-Japanese (kango): two Chinese characters read with their on'yomi, readings borrowed from Chinese. They are \emph{nichi}, \textquotedblleft{}day, sun\textquotedblright{} (Middle Chinese \emph{nyit}) and \emph{jō}, \textquotedblleft{}usual, regular\textquotedblright{} (\emph{dzyang}) (KANJIDIC2; Middle Chinese in Baxter's transcription).
\end{cousinweb}

\begin{grammarlens}[title={What you've built}]
One more word for everyday life at A2.
\end{grammarlens}

\subsection*{Your turn}
\begin{itemize}
\raggedright
  \item \emph{Say it:} \emph{nichijō}
  \item \emph{Say it:} \emph{nichijō}, once more
  \item \emph{Say it:} \emph{nichijō wa tanoshii desu} — Everyday life is enjoyable.
  \item \emph{Recall:} say the Japanese for usually, then the Japanese for a habit, then say \emph{nichijō} again
\end{itemize}

\begin{tcolorbox}[breakable,colback=teal!4,colframe=teal!35!black,title={Before you move on}]
What does \ja{にちじょう} mean? (\textquotedblleft{}Everyday life; daily\textquotedblright{}.) Say it once more.
\end{tcolorbox}
\section[\texorpdfstring{\ja{くらす} (kurasu)}{くらす (kurasu)}]{\ja{くらす} (kurasu) — to get by; to live one's days}

\label{lesson:JA-C154-kurasu}



\begin{quote}
Before the new one: say the Japanese for a habit, then the Japanese for everyday life.
\end{quote}

\subsection*{You'll want to know: \ja{くらす}}
\textbf{\ja{くらす}} — \emph{kurasu} — \textquotedblleft{}to get by; to live one's days\textquotedblright{}.

\begin{quote}
\textbf{\ja{いなかで} \ja{くらします。}} — \emph{inaka de kurashimasu} — \textquotedblleft{}I live in the countryside.\textquotedblright{}
\end{quote}

Polite \textbf{\ja{くらします}}, \emph{kurashimasu}; a godan verb.

\begin{cousinweb}
\textbf{\ja{くらす}} (\emph{kurasu}) is a native Japanese word (wago), inherited rather than borrowed from Chinese; its kanji was attached later as a kun'yomi, and its older origin is not settled (JMdict, KANJIDIC2).
\end{cousinweb}

\begin{grammarlens}[title={What you've built}]
One more word for everyday life at A2.
\end{grammarlens}

\subsection*{Your turn}
\begin{itemize}
\raggedright
  \item \emph{Say it:} \emph{kurasu}
  \item \emph{Say it:} \emph{kurasu}, once more
  \item \emph{Say it:} \emph{inaka de kurashimasu} — I live in the countryside.
  \item \emph{Recall:} say the Japanese for a habit, then the Japanese for everyday life, then say \emph{kurasu} again
\end{itemize}

\begin{tcolorbox}[breakable,colback=teal!4,colframe=teal!35!black,title={Before you move on}]
What does \ja{くらす} mean? (\textquotedblleft{}To get by; to live one's days\textquotedblright{}.) Say it once more.
\end{tcolorbox}
\section[\texorpdfstring{\ja{なれる} (nareru)}{なれる (nareru)}]{\ja{なれる} (nareru) — to get used to}

\label{lesson:JA-C154-nareru}



\begin{quote}
Before the new one: say the Japanese for everyday life, then the Japanese for to get by.
\end{quote}

\subsection*{You'll want to know: \ja{なれる}}
\textbf{\ja{なれる}} — \emph{nareru} — \textquotedblleft{}to get used to\textquotedblright{}.

\begin{quote}
\textbf{\ja{こどもは} \ja{がっこうに} \ja{なれます。}} — \emph{kodomo wa gakkō ni naremasu} — \textquotedblleft{}Children get used to school.\textquotedblright{}
\end{quote}

Polite \textbf{\ja{なれます}}, \emph{naremasu}; an ichidan verb.

\begin{cousinweb}
\textbf{\ja{なれる}} (\emph{nareru}) is a native Japanese word (wago), inherited rather than borrowed from Chinese; its kanji was attached later as a kun'yomi, and its older origin is not settled (JMdict, KANJIDIC2).
\end{cousinweb}

\begin{grammarlens}[title={What you've built}]
One more word for everyday life at A2.
\end{grammarlens}

\subsection*{Your turn}
\begin{itemize}
\raggedright
  \item \emph{Say it:} \emph{nareru}
  \item \emph{Say it:} \emph{nareru}, once more
  \item \emph{Say it:} \emph{kodomo wa gakkō ni naremasu} — Children get used to school.
  \item \emph{Recall:} say the Japanese for everyday life, then the Japanese for to get by, then say \emph{nareru} again
\end{itemize}

\begin{tcolorbox}[breakable,colback=teal!4,colframe=teal!35!black,title={Before you move on}]
What does \ja{なれる} mean? (\textquotedblleft{}To get used to\textquotedblright{}.) Say it once more.
\end{tcolorbox}
